% SOURCE_AUTHORITY: sga5_fr_workpass.tex, lines 6602--6633 (LF-normalized unit).
% SOURCE_SHA256: 791F4EFFC5E02832D5D77ED03518C8156D6F07E4C8238B03545DB93D883FBB28
\subsection*{6.24}
Se $\Lambda$ for anulada por uma potência de um número primo $\ell\ne p$, o segundo membro de 6.23.1 poderá ser reescrito na forma
\[
\sum_{g\in G_{\natural}}
\Tr_{\mathbb Z_{\ell}[G]}(f_*\underline c_{\ell,U,X})(g^{-1})
\Tr_{\Lambda}(gv_0),
\]
com as notações de 6.17. A fórmula 6.17.1 mostra que os termos locais definidos por Grothendieck em (XII 6.2) coincidem com os termos locais de Verdier, pelo menos quando o endomorfismo $\varphi$ de $G$ considerado em (loc. cit.) é a identidade; mas, como já se observou, essa restrição não é grave. Se $Y$ é próprio sobre $S$, combinando 6.23.1 com a fórmula de Lefschetz para $Y\to S$ (III 4.7), obtém-se uma fórmula que generaliza (XII 6.3) (e, com maior razão, a fórmula de Lefschetz para Frobenius demonstrada em (SGA $4^{1/2}$ Rapport) e a parte I da presente exposição).

\section*{Bibliografia}

\begin{enumerate}[label={[\arabic*]}]
\item Cartan, H. e Eilenberg, S. -- \textit{Homological Algebra}, Princeton University Press, 1956.
\item Grothendieck, A. -- Formule de Lefschetz et rationalité des fonctions $L$, Séminaire Bourbaki 64/65, n\textsuperscript{o} 279, em \textit{Dix exposés sur la cohomologie des schémas}, North Holland, Pub. Co., 1968.
\item Illusie, L. -- Complexe cotangent et déformations I, II, Lecture Notes in Mathematics 239, 283, Springer-Verlag, 1971, 1972.
\item Langlands, R. P. -- Modular forms and $\ell$-adic representations, em \textit{Modular Functions of One Variable II}, Lecture Notes in Mathematics 349, Springer-Verlag, 1973.
\item Stallings, J. -- Centerless groups -- an algebraic formulation of Gottlieb's theorem, \textit{Topology} 4, p. 129--134, 1965.
\item Verdier, J.-L. -- The Lefschetz fixed point formula in étale cohomology, em \textit{Proceedings of a conference on local fields}, Springer-Verlag, 1967.
\end{enumerate}

SGA $4^{1/2}$: \textit{Cohomologie étale}, por P. Deligne, Lecture Notes in Mathematics 569, Springer-Verlag, 1977.

Em SGA $4^{1/2}$:
\[
\text{Th. Finitude: Teoremas de finitude em cohomologia }\ell\text{-adique;}
\]
\[
\text{Ciclo: A classe de cohomologia associada a um ciclo;}
\]
\[
\text{Relatório: Relatório sobre a fórmula dos traços.}
\]
